\documentclass{article}

\usepackage{amsmath}
\usepackage[boxed,vlined,boxruled]{algorithm2e}


\title{Divide and Conquer}
\author{Kanishk Goel}
\date{22 February, 2026}
\parindent 0px

\SetKwInput{KwAssume}{Simplifying assumption}
\SetAlgoCaptionLayout{centerline}

\begin{document}
	\SetAlgoCaptionSeparator{}
	\renewcommand{\thealgocf}{} 
	\renewcommand{\AlCapNameFnt}{\bfseries}
	\renewcommand{\AlCapFnt}{}
	\DontPrintSemicolon
	\renewcommand{\algorithmcfname}{}
	\maketitle
	\SetAlgoNoEnd
	\section{MergeSort}

		
	\begin{algorithm}[H]
		\SetAlgoLined
		\caption{Problem: Sorting}
		\KwIn{An array of $n$ numbers in arbitrary order.}
		\KwOut{Array with the same numbers, sorted in ascending order}
	\end{algorithm}
	\vspace{5mm}
	
	\begin{algorithm}[H]
		\SetAlgoLined
		\caption{MergeSort}
		\KwIn{Array $A$ of $n$ distinct integers.}
		\KwOut{Array with the same integers, sorted in ascending order}
		\BlankLine
		\If{$n \leq 1$}{
			\Return{$A$}\
		}\
		
		$C :=$ recursively sort first half of $A$\;
		
		$D :=$ recursively sort second half of $A$\;
		
		\Return{$\textnormal{Merge}(C,D)$}\;
		  
	\end{algorithm}
	\vspace{5mm}
	
\begin{algorithm}[H]
	\SetAlgoLined
	\caption{Merge}
	\KwIn{Sorted arrays $C$ and $D$ (length $n/2$ each)}
	\KwOut{Sorted array $B$ (length $n$)}
	\KwAssume{$n$ is even}
	\BlankLine
	
	$i := 1$\;
	$j := 1$\;
	
	\For{$k := 1$ \KwTo $n$}{
		\If{$i \le n/2$ \textbf{and} ($j > n/2$ \textbf{or} $C[i] \le D[j]$)}{
			$B[k] := C[i]$\;
			$i := i + 1$\;
		}
		\Else{
			$B[k] := D[j]$\;
			$j := j + 1$\;
		}
	}
\end{algorithm}
\vspace{5mm}

	
\begin{algorithm}[H]
	\SetAlgoLined
	\caption{Problem: Counting Inversions}
	\KwIn{An array $A$ of distinct integers.}
	\KwOut{The number of inversions of $A$ i.e pairs $(i,j)$ of array indices with $i \le j$ and $A[i] \ge A[j]$.}
\end{algorithm}
\vspace{5mm}

Used in Collaborative Filtering which is a technique for generating recommendations.

\begin{algorithm}[H]
	\SetAlgoLined
	\caption{CountInversions}
	\KwIn{array $A$ of $n$ distinct integers.}
	\KwOut{sorted array $B$ with the same integers, and the number of inversions of $A$.}
	\BlankLine
	\If{$n=0 \lor n= 1$}{
		\Return{$(A,0)$}\;
	}\;
	\Else{
		$(C, leftInv)$ := CountInversions$($first half of A$)$ \;
		$(D, rightInv)$ := CountInversions$($second half of A$)$ \;
		$(B, splitInv)$ := CountSplitInversions$(C,D)$ \;
		\Return{$(B, leftInv+rightInv+splitInv)$}
	}\;
\end{algorithm}
\vspace{5mm}

\begin{algorithm}
	\SetAlgoLined
	\caption{CountSplitInversions}
	\KwIn{sorted arrays $C$ and $D$ (length $n/2$ each).}
	\KwOut{sorted array $B$ (length $n$).}
	\KwAssume{$n$ is even}
	\BlankLine
	$i$ := $1$, $j$ := $1$, $splitInv$ := $0$ \;
	\BlankLine
	\For{$k$ := $1$ to $n$}{
		\If{$C[i] < D[j]$}{
			$B[k]$ := $C[i]$, $i$ := $i+1$
		}\;
		\Else{
			$B[k]$ := $D[j]$, $j$ := $j + 1$ \;
			$splitInv$ := $splitInv + (\frac{n}{2} - i + 1)$
		}\;
	}\;
	\Return{$(B,splitInv)$}
\end{algorithm}

The idea used is equivalent to merge subroutine of mergesort.

\section{QuickSort}
The pseudocode for quicksort: 

\begin{algorithm}[H]
	\SetAlgoLined
	\caption{QuickSort}
	\KwIn{array $A$ of n distinct integers, left and right endpoints $l,r \in \{1,2,3 \ldots ,n\}.$}
	\KwOut{Elements of the subarray $A[l],A[l+1], \ldots , A[r]$ are sorted from smallest to largest.}
	\BlankLine
	\If{$l \geq r$}{
		\Return{}
	}
	$i \leftarrow$ ChoosePivot$(A,l,r)$ \;
	swap $A[l]$ and $A[i]$ \;
	$j \leftarrow$ Partition$(A,l,r)$ \;
	QuickSort$(A,l,j-1)$ \;
	QuickSort$(A,j+1,r)$ \;
\end{algorithm}	
\vspace{5mm}
\begin{algorithm}[H]
	\SetAlgoLined
	\caption{Partition}
	\KwIn{array $A$ of $n$ distinct integers, left and right endpoints $l,r \in \{1,2 \ldots , n\} with l \leq r$.}
	\KwOut{Elements of the subarray $A[l], \ldots , A[r]$ are partitioned around $A[l]$. Final output is the position of the pivot element.}
	\BlankLine
	$p \leftarrow A[l]$ \;
	$i \leftarrow l+1$ \;
	\For{$j \leftarrow l+1$ to $r$}{
		\If{$A[j] \ge p$}{
			swap $A[j]$ and $A[i]$ \;
			$i \leftarrow i +1$ \;
		}
	}
	swap $A[l]$ and $A[i-1]$ \;
	\Return{$i-1$}
\end{algorithm}
\vspace{2mm}	
Now the idea of choosing pivot is important, randomized pivot selection has an average running time of $O(n \log n)$

\section{Selection Problem}
\begin{algorithm}[H]
	\SetAlgoLined
	\caption{Problem: Selection}
	\KwIn{An array $A$ of $n$ numbers, in arbitrary order, and an integer $i \in \{1,2, \ldots , n\}$. }
	\KwOut{The $i^{th}$ smallest element of $A$ or the $i^{th}$ order statistic.}
	\BlankLine
\end{algorithm}	
\vspace{2mm}
There are two algorithms one is RSelect which is a randomized select algorithm with an average case of $O(n)$. The other is DSelect.

\begin{algorithm}[H]
	\SetAlgoLined
	\caption{RSelect}
	\KwIn{array $A$ of $n \geq 1$ distinct numbers, and an integer $i \in \{1,2, \ldots , n\}$}
	\KwOut{the $i^{th}$ order statistic of $A$.}
	\BlankLine
	\If{$n=1$}{
		\Return{$A[1]$}
	}
	choose pivot element $p$ at random from $A$ \;
	partition $A$ around $p$ \;
	$j \leftarrow p$'s position in partitioned array \;
	\If{$j = i$}{
		\Return{$p$}
	} \;
	\ElseIf{$j > i$}{
		\Return{RSelect(first part of $A$, $i$)}
	} \;
	\Else{
		\Return{RSelect(second part of $A$, $i$)}
	}
\end{algorithm}	
\vspace{5mm}
\begin{algorithm}[H]
	\SetAlgoLined
	\caption{DSelect}
	\KwIn{array $A$ of $n \geq 1$ distinct numbers, and an integer $i \in \{1,2, \ldots , n\}$}
	\KwOut{the $i^{th}$ order statistic of $A$.}
	\BlankLine
	\If{$n=1$}{
		\Return{$A[1]$}
	} \;
	\For{$h \leftarrow 1$ to $\frac{n}{5}$}{
		$C[h] \leftarrow$ middle element from the $h^{th}$ group of 5
	} \;
	$p \leftarrow DSelect(c,\frac{n}{10})$\;
	partition $A$ around $p$ \;
	$j \leftarrow p$'s position in partitioned array \;
	\If{$j = i$}{
		\Return{$p$}
	} \;
	\ElseIf{$j > i$}{
		\Return{DSelect(first part of $A$, $i$)}
	} \;
	\Else{
		\Return{DSelect(second part of $A$, $i$)}
	}
\end{algorithm}



\end{document}